\subsection{Expectation bounds}
\label{sec:expectation_bounds}

Our goal is to find a bound on the expectation of certain random variables
which are products of $\bar{B}_{ef}$ along a circuit.
To do this, we will need to bound the probabilities of different
subgraphs when exploring $G$.
This requires us to introduce the concept of consistent
edges and their multiplicity.

\begin{definition}\label{def:incon_edges}
Let $\gamma = (e_1, \ldots, e_{2k})$ be a sequence of half edges
of even length, with 
$\vec{E}(\gamma)$ its set of half edges,
and $E(\gamma) = \{ \{e_{2i-1}, e_{2i}\} \mbox{ for $i\in [k]$} \}$ its set of edges (unordered pairs, thus undirected).
\begin{itemize}
    \item The {\em multiplicity of a half edge} $e \in \vec{E}(\gamma)$
    is $m_\gamma(e) = \sum_{t=1}^{2k} 1_{\{ e_t = e \}}$.
    \item The {\em multiplicity of an edge} $\{h_1, h_2\} \in E(\gamma)$,
    is 
    $m_\gamma(\{h_1,h_2\}) = 
    \sum_{t=1}^k 1_{\{ \{e_{2t-1}, e_{2t}\} = \{h_1, h_2\} \}}$.
    \item An edge $\{h_1, h_2\}$ is {\em consistent} if 
    $m_\gamma(h_1) = m_\gamma(h_2) = m_\gamma(\{h_1, h_2\})$.
\end{itemize}
\end{definition}

\begin{lemma}
\label{expectation_bound}
Let  
$\gamma = (e_1, \ldots, e_{2k})$ be a sequence of half edges
of even length, with
$M$ and $\bar{M}$ the matching matrix and its
centered version generated by a uniform matching in the 
configuration model. 
Then for $1 \leq k \leq \sqrt{|E|}$ and $0 \leq t_0 \leq k$
we have that
\[
\left| 
\E 
\prod_{t=1}^{t_0} \bar{M}_{e_{2t - 1} e_{2t}}
\prod_{t=t_0+1}^{k} M_{e_{2t - 1} e_{2t}}
\right|
\leq 
C \cdot 2^b \cdot 
\left( \frac{1}{|E|} \right)^\nedge
\left( \frac{3 k}{\sqrt{|E|}} \right)^{\nedge_1}
\]
where
$b = $ number of inconsistent edges of multiplicity one occuring
before $t_0$,
$\nedge_1 = $ number of consistent edges with multiplicity one
occuring before $t_0$,
$\nedge = |E(\gamma)|$,
and $C$ is a universal constant.
\end{lemma}

\begin{proof}
Recall the form of the matrices
\begin{equation*}
    M = 
    \left(
    \begin{array}{cc} 
        0 & M_1 \\
        M_1^* & 0 
    \end{array}
    \right )
    \quad \mbox{and} \quad
    \bar{M} = M - 
    \frac{1}{|E|} \left(
    \begin{array}{cc} 
        0 & \1 \1^* \\
        \1 \1^* & 0 
    \end{array}
    \right ).
\end{equation*}
Matrix $M_1 \in \R^{|E| \times |E|}$ is a random permutation matrix between 
$n d_1 = |E|$ and $m d_2 = |E|$ half edges.
Therefore, $M_1$ is distributed exactly the same as a matching matrix
of a random $|E|$-lift of a single edge,
and the same holds for its centered version $M_1 - \frac{1}{|E|} \1 \1^*$.
The only paths $\gamma$ that contribute in this bipartite setting must alternate
between the bipartite sets and avoid the 0 blocks,
otherwise the bound holds trivially.
For one of these paths $\gamma$
assume, without loss of generality, that the path starts in set $V_1$.
Then define the transformed path 
$\gamma' = (e_1', \ldots, e_{2k}') = (e_1, e_2, e_4, e_3, e_5, \ldots)$,
i.e.\ with every other pair in $\gamma$ in reverse order.
Note that
\begin{equation}
\prod_{t=1}^{t_0} \bar{M}_{e_{2t - 1} e_{2t}}
\prod_{t=t_0+1}^{k} M_{e_{2t - 1} e_{2t}}
=
\prod_{t=1}^{t_0} (\bar{M}_1)_{e'_{2t - 1} e'_{2t}}
\prod_{t=t_0+1}^{k} (M_1)_{e'_{2t - 1} e'_{2t}} .
\end{equation}
Then the Lemma holds by
\cite{bordenave2015}, Proposition 28.
\end{proof}

% The next Lemma follows from \citet*{mckay1981a}, Theorem 3.5. 
% We use the following notation: 
%  Let $H$ be a subgraph of $G$ with vertex set $\{v_1,v_2,\dots,v_k\}$. 
% Let $d_i$ be the degree of $v_i$ in $G$, 
% so $d_i=d_1$ if $v_i$ is in the set $V_1$ and $d_2$ if not. 
% Also, denote by $h_i$ the degree of $v_i$ in $H$. 
% For natural numbers $x$ and $t$, we use
% \[
% (x)_{t} = x (x-1) (x-2) \cdots (x-t+1)
% \]
% to denote the falling factorial. 
% \begin{lemma}
% \label{lem:McKay}
% Let $H\subset K_{n,m}$ such that 
% $|E(H)|=o(n)$ and $G \sim \Gg$.
% Then
% \[
% \P(H\subset G) 
% \leq 
% \frac{ \prod_{i=1}^k (d_i)_{h_i} }{ (n d_1 - 4 d^2)_{ |E(H)| } } ,
% \]
% and
% \[
% \P(H\subset G)
% \geq 
% \frac{ \prod_{i=1}^k (d_i)_{h_i} }{ (n d_1 - 1)_{|E(H)|} }
% \left( 
%   \frac{ n d_1 - |E(H)| - 5 d^2 }{ n d_1 - c |E(H)| - 5 d^2}
% \right)^{ |E(H)| } ,
% \]
% for some explicit constant $c < 1$, where $d = \max (d_1,d_2)$. 
% \end{lemma}

% Crucially, we will use Lemma~\ref{lem:McKay}
% to show that the appearance of edges in
% random bipartite biregular graphs are weakly correlated, 
% as long as the number of edges is not too large. 
% Computations are similar to those carried out in \citet*{brito2016}. 

% \begin{lemma}\label{one_edge}
% Let $G\sim \mathcal{G}(n,m,d_1,d_2)$ and let $H\subset K_{n,m}$ 
% such that $|E(H)| = o(n)$. 
% Let $e$ be an edge not in $H$, such that 
% $H$ and $H\cup\{e\}$ have the same number of connected components,
% and let $n_1 = m$, $n_2 = n$.
% \begin{enumerate}[(i)]
% \item If $e$ and $H$ share an endpoint of degree $d_i$ in $G$, it holds that
% \[
% \P(e \in G | H \subset G) \leq \frac{d_i-1}{n_i} + O\left( \frac{|E(H)|}{n^2} \right) . 
% \]
% \item If $e$ has exactly one endpoint in $H$, 
% and this vertex has degree one in $H$
% and degree $d_i$ in $G$, then
% \[
% \P(e \in G | H \subset G)= \frac{d_i-1}{n_i} + O\left( \frac{|E(H)|}{n^2} \right) .
% \]
% \item If $e$ has endpoints in both $V_1$ and $V_2$, then we can use the bound in (i)
% with $i = 1$ or $2$.
% \end{enumerate}
% \end{lemma}

% \begin{proof}
% Let $e = (u, v)$,  where $v$ is the shared endpoint in $H$.
% Assume $v$ is degree $d_2$ in $G$; the results for $d_1$ are analogous.
% Vertex $u \not \in V(H)$ must then have degree $d_1$ in $G$.
% For part (i), notice that the graph $H \cup \{ e \}$ 
% has at least one vertex $v$ with degree equal to $h_v + 1$.
% We use the upper bound in Lemma \ref{lem:McKay}
% and the geometric series to obtain
% \begin{align*}
% \P(H\cup\{e\} \subset G) 
% & \leq 
% \frac{ d_1 (d_2 - h_v) }{ n d_1 - 4d^2 - |E(H)| } \P( H \subset G) \\
% & \leq
% \frac{ d_2 - 1 }{ n } 
% \left( 1 + O\left( \frac{ |E(H)| }{ n } \right) \right) 
% \P( H \subset G).
% \end{align*}
% For part (ii), note that $h_v=1$ and employ the lower bound in Lemma~\ref{lem:McKay} to get equality. 
% \end{proof}

% \begin{remark}
% Assuming $m \geq n$, and using $\frac{d_1 - 1}{m} = \frac{d_2 - 1}{n} + \frac{1}{n} - \frac{1}{m}$,
% we see that the bound with $d_1$ is weaker by a factor of $O(\frac{1}{n})$.
% Finally, $\P (e \in G | H) \leq \frac{d_1}{m} = \frac{d_2}{n}$ 
% for any isolated edge.
% \end{remark}
% \begin{remark}%[Probability of tree edges]
% The centering matrix $S$ has entries $(d_2 -1)/n$ and $(d_1-1)/m$.
% During the exploration process,
% Lemma~\ref{one_edge} states that the probability of adding a new edge
% differs from the entries of the centering matrix
% (which are close to but not exactly the expectation)
% by an order $1/n^2$ correction
% {\it precisely} when adding that edge creates a two-path.
% See Figure~\ref{fig:tree_edges}.
% \end{remark}

% \begin{figure}[t]
% \centering
% \includegraphics[width=0.5\linewidth]{tree_vs_non_tree_edges.png}
% \caption{The different cases of Lemma~\ref{one_edge},
% which concerns $P(e \in G | H)$:
% the probability of an edge $e$ existing in the random graph $G$, 
% given an observed subgraph $H$, shown in green.
% The new edge $e$ connects at vertex $v$.
% All the green vertices lie in $H$.
% When the connecting vertex has degree one in
% subgraph $H$, i.e.\ $h_v = 1$ and case (ii),
% it induces a two-path. 
% The probability of that edge is closer to 
% the entries of the centering matrix $S$ 
% if $h_v > 1$, case (i).
% }
% \label{fig:tree_edges}
% \end{figure}

% Let $\gamma=(\gamma_1,\gamma_2,\cdots,\gamma_{2k})$ 
% be a circuit obtained by the concatenation of $2k$ non-backtracking walks of length $\ell$. 
% These circuits will appear when we apply the trace method in Section~\ref{sec:norm_bounds}.
% The circuit $\gamma$ is traversed in this order: 
% We start at the initial vertex of $\gamma_1$, move along this path until we meet $\gamma_2$,
% continue along this path, etc. 
% Denote by 
% $\vec{E}(\gamma)$ 
% the set of oriented edges traversed by $\gamma$,
% and $E(\gamma)$ the same set of edges without orientation. 
% A subpath of $\gamma$ is just an ordered path of edges traversed as described above. 
% Define
% \begin{align}
% \label{def:X_gamma}
% X_{\gamma} = \prod_{ef \in \gamma} \left( \bar{B}_{ef} \right)^{ m_{ef} },
% \end{align}
% where $e, f\in \vec{E}(\gamma)$,
% and $ef \in \gamma$ means that the oriented path $ef$ 
% is a subpath of $\gamma$ when traversed as described above,
% and $m_{ef}$ is the total number of times we traverse $ef$ in $\gamma$, i.e.\
% its {\it multiplicity}. 
% The notation $e^{-1}$ refers to the edge $e$ traversed in the reversed orientation.
% The main result of this section is the following Theorem:
% \begin{theorem}
% \label{expectation_bound}
% Let $\gamma=(\gamma_1,\gamma_2,\cdots,\gamma_{2k})$ 
% be a circuit obtained by the concatenation of $2k$ non-backtracking walks of length $\ell$. 
% Suppose $\gamma$ visits $\nedge=|E(\gamma)|=o(n)$ different edges and 
% $\nvright = |V(\gamma) \cap V_2|$ vertices in the right set.
% Then we have the bound
% \begin{align}
% \label{exp_bound_ineq}
% |\mathbb{E} \left( X_{\gamma} \right) |
% \leq 
% \frac{ ((d_1 - 1) (d_2 - 1))^{\nedge/2} }{ m^\nvright n^{\nedge - \nvright} }  
% \left( \frac{d_1 - 1}{d_2 - 1} \right)^{\psi}
% %\left( d-1 \right)^{\nedge - 2 k \ell}
% %\frac{ (d_1-1)^{\nvright}(d_2-1)^{\nedge - \nvright} }{ m^\nvright n^{\nedge - \nvright} }  
% %\left( \frac{d_2}{d_2 - 1} \right)^s
% \left( \frac{c \nedge}{n} \right)^{ \omega }
% (1 + o(1)),
% \end{align}
% where 
% \[
% \omega 
% = 
% \left\lceil 
% \frac{
% \sum_{e f \in \gamma} 
% \mathbf{1} \{m_{e f} + m_{f^{-1} e^{-1}} = 1 \} 
% }
% {6 d^2}
% \right\rceil,
% \quad
% \psi = \nvright - \nedge/2,
% \]
% $d = \max(d_1, d_2)$,
% and 
% %$\max(1, \omega/2) \leq s \leq 6d^2 \omega$.
% %\todo{how big can $s$ be?}
% $c = \left( \frac{d_2}{d_2-1} \right)^{6d^2}$.
% \end{theorem}
% % OLD VERSION
% % \begin{theorem}
% % \label{expectation_bound}
% % Let $\gamma=(\gamma_1,\gamma_2,\cdots,\gamma_{2k})$ 
% % be a circuit obtained by the concatenation of $2k$ non-backtracking walks of length $\ell$. 
% % Suppose $\gamma$ visits $\nedge=|E(\gamma)|=o(n)$ different edges and 
% % $\nvright = |V(\gamma) \cap V_2|$ vertices in the right set.
% % Let $d = \max(d_1, d_2)$ and
% % \[
% % \omega 
% % = 
% % \left\lceil 
% % \frac{\sum_{ef \in \gamma} \mathbf{1} \{m_{e f}+m_{f^{-1} e^{-1}}=1\} }{6 d^2}
% % \right\rceil.
% % \]
% % Depending on the circuit $\gamma$, we obtain two cases:
% % \begin{enumerate}[(a)]
% % \item
% % Every two-path is visited at least twice, and
% % $\omega = 0$:
% % \begin{align}
% % \label{exp_bound_ineq}
% % \mathbb{E} \left( X_{\gamma} \right)
% % \leq 
% % \frac{ (d_1-1)^{\nvright}(d_2-1)^{\nedge - \nvright} }{ m^\nvright n^{\nedge - \nvright} }  
% % \left( \frac{d_2}{d_2 - 1} \right)
% % (1 + o(1)).
% % \end{align}
% % \item 
% % Some two-paths are visited only once, and
% % $\omega > 0$:
% % \begin{align}
% % \label{exp_bound_ineq_b}
% % \mathbb{E} \left( X_{\gamma} \right)
% % \leq 
% % \frac{ d^\nedge }{ m^\nvright n^{\nedge - \nvright} } 
% % \left( \frac{\nedge}{n} \right)^{ \omega } (1 + o(1)).
% % \end{align}
% % \end{enumerate}
% % \end{theorem}

% We now outline the steps in the proof of Theorem~\ref{expectation_bound}
% and gather some useful Lemmas before carrying out the proof.
% We choose an ordering of the set of {\it undirected} edges visited by $\gamma$.
% With this ordering, we define a filtration, or sequence of nested sigma algebras,
% $\{\mathcal{F}_t \}_{t=1}^\nedge$ 
% each containing the information of the first $t$ edges in this order. 
% We use the tower of expectation to bound this quantity.
% At each step, a new edge is removed from the filtration and 
% we are able to improve our current bound via Lemma~\ref{one_edge}. 
% The ordering of the edges is done in a way that allows us to use 
% part (ii) of Lemma~\ref{one_edge} a maximal number of times,
% which will add $\omega$ extra factors of $\nedge/n$.
% These extra factors come from ``good'' edges,
% appearing in two-paths which are traversed exactly once.

% We start by describing the ordering of the edges. 
% Let $E = \{e_i\}_{i=1}^\nedge$ be a set of undirected edges. 
% A permutation $\pi$ of $[\nedge]$
% can be identified with an ordering in $E$ by taking the first edge to be $e_{\pi(1)}$, 
% the second to be $e_{\pi(2)}$, etc. 
% For a subset $F\subseteq  E$ of edges, define
% \[
% N(F)=\{ e\in E:~\mbox{$e$ shares a vertex with some $f\in F$} \}
% \]
% of \textit{neighbors} of $F$. 
% Notice that the orientation is not relevant in this definition. 

% Given an ordering $\pi$, we say that the edge $e_{\pi(j)}$ is a
% {\it good edge} if there is exactly one value of $i\leq j-1$ 
% such that $e_{\pi(i)}\in N(e_{\pi(j)})$.
% In other words, 
% the $j$th edge is good if it has exactly one neighbor among the previous edges.
% Thus, when we add a good edge to the graph induced by the previous edges, 
% it must induce a two-path. 
% The following Lemma tells us how many good edges we can guarantee with our 
% ordering $\pi$, and gives a recipe to construct that ordering.

% \begin{lemma}
% \label{lem:good_ordering}
% Let $F$ be some set of undirected
% edges in a graph with maximal degree $d$,
% with $p_2(F)$ equal to the number
% of two-paths in $F$.
% Then there exists an ordering $\pi$ of the edges with
% at least
% $
% \left\lceil \frac{p_2(F)}{6 d^2} \right\rceil
% $
% good edges.
% \end{lemma}
% \begin{proof}
% We now construct such an ordering by working with three sets of edges
% $E^1$, $E^2$, and $E^3$.
% At time $t=0$, start with the sets $E^1_0=\emptyset$, $E^2_0=\emptyset$ and $E^3_0 = F$. 
% For times $t \geq 1$, we build the ordering $\pi$ iteratively:
% \begin{enumerate}
% \item Choose $e_i,e_j\in E^3_{t-1}$ such that $e_i \in N(e_j)$, 
% 	i.e.\ $e_i e_j$ is a two-path. 
% \item Set $\pi(2t-1)=i$ and $\pi(2t)=j$. 
% \item Set $E^1_t = E^1_{t-1}\cup\{e_i,e_j\}$, 
% $E^2_t = E^2_{t-1}\cup N(\{e_i,e_j\})\backslash E^1_t$ and 
% $E^3_t = F\backslash(E^1_t\cup E^2_t)$.
% \end{enumerate}

% At all times, $E^1_t, E^2_t$ and $E^3_t$ form a disjoint partition of $F$. 
% They correspond to:
% the set of edges already ordered $E^1$, 
% the set of the neighbors of the edges already ordered $E^2$,
% and the complement of those two sets $E^3$. 
% Also, it is not hard to check that $e_{2t}$ is a good edge for all $t$. 
% This process will end at step $t=T$ 
% when one of the following exclusive events happen: 
% \begin{enumerate}[(a)]
% \item The set $E^3_T=\emptyset$ and we defined $\pi(t)$ for 
% all $1 \leq t \leq 2T$. 
% We now arbitrarily set $\pi(t)$ for times $2T+1\leq t\leq |F|$ 
% using the edges in $E^2_T$. 
% By construction, we have at least $T$ good edges. 
% Since at each step we remove from $E^3$ at most $3d$ edges,
% and each edge participates in less than $2d$ two-paths,
% we conclude that
% \[
% T \geq \left\lceil \frac{p_2(F)}{6d^2} \right\rceil .
% \]
% \textcolor{magenta}{[Note: this adds $s = T$ isolated edges,
% since edges in $E_T^2$ are connected.]}
% \item After $T$ steps, no two edges in $E^3_T$ form a two-path. 
% By construction, $E^1_T$ and $E^3_T$ are disconnected. 
% We use this to continue the ordering for times $2T+1\leq t\leq |F|$:
% Because $\gamma$ is a connected walk, 
% for each $e_i \in E^3_t$ there exists some $e_j \in E^2_t$ 
% such that $e_i e_j$ or $e_j e_i$ is a subpath of $\gamma$. 
% Set ${\pi(2 t+1) } = j$ and ${\pi(2t+2)}=i$. 
% With this choice, $e_i$ is good. 
% We update $E^1_{t+1}$, $E^2_{t+1}$, and $E^3_{t+1}$ as before.
% The process is repeated until time $t=T^*$, when $E^3_{T^*}=\emptyset$. 
% The number of good edges is again at least 
% \[
% T^* \geq \left\lceil \frac{p_2(F)}{6d^2} \right\rceil .
% \]
% \textcolor{magenta}{[Note: this adds
% $s = T \leq T^*$
% isolated edges, none after $T$ steps.]}
% \end{enumerate}
% \end{proof}
% \textcolor{magenta}{Above considerations mean
% we need to upper bound the number of good edges as well...}

% \begin{remark}
% The above Lemma gives an algorithm for ordering some set of edges $F \subseteq E$
% in a way that guarantees good edges.
% Suppose this process finishes at time $T$.
% The edges in $E_T^1$ are added in a way that they form a disjoint
% set of two-paths, and half of these edges are good.
% We note that we are free to add the edges left in $E \setminus E_T^1$ 
% in any way without decreasing the number of good edges.
% \end{remark}

% We will also need the following Lemma to bound the expectation
% of a nearly centered Bernoulli random variable.
% \begin{lemma}[\citet*{brito2015}, Lemma 11]
% \label{bernoulli_bound}
% Let $0 \leq p,q \leq 1$ and $r$ such that $q \leq p+r$. 
% Let $X\sim \mathrm{Bernoulli}(q)$. 
% Then for any integer $m>1$, the expectation
% $
% \mathbb{E} \left( |X - p|^m \right)\leq p+r.
% $
% \end{lemma}
% \begin{proof}
% Assume $q<p$. Then,
% \[
% \mathbb{E}(|X-p|^m)\leq (1-p)^mp+p^m\leq p .
% \]
% The latter inequality follows easily by noting that it is satisfied
% for $m=2$, and that $(1-p)^m p +p^m$ is a decreasing function of $m$
% for all $0 \leq p\leq 1$.

% If $q>p$, write $q=p+r'$ with $0<r'<r$. 
% We get:
% \[
% \mathbb{E}(|X-p|^m)\leq(1-p)^m(p+r')+p^m\leq (1-p)^m p+p^m+r'\leq p+r~,
% \]
% due to similar considerations.
% \end{proof}

% \todo{I believe we will need a Lemma like this for the case of $m=1$:}\textcolor{green}{Alternatively, we can include it as a second statement in Lemma 11.}
% \begin{lemma}
% For $p,q,r$ as before, we have that
% $ \E |X - p| \leq 2p + r$.
% \end{lemma}

% \begin{proof}[Proof of Theorem~\ref{expectation_bound}] 

% First, consider the following elementary chain of inequalities:
% \begin{equation}
% \E ( X_\gamma )
% \leq 
% \left| \E ( X_\gamma ) \right|
% \leq
% \E \left(
%     | X_\gamma |
% \right)
% =
% \E
% \left(
% \prod_{ef \in \gamma}
% \left| \bar{B}_{ef} \right|^{m_{ef}} 
% \right) 
% ,
% \end{equation}
% where the second inequality is a consequence of Jensen.
% We will work with the unsigned matrix $|\bar B|$
% whose entries are just the absolute values of the entries in $\bar B$.
% \todo{The above inequalities are correct, 
% but it is possible that they are too weak. 
% We really just need to bound $| \E (X_\gamma) |$.}
% % \begin{align*}
% % \E ( X_\gamma )
% % &\leq 
% % \left| \E ( X_\gamma ) \right| \\
% % &=
% % \left|
% % \E \left(
% %   \prod_{ef \in \gamma} \bar{B}_{e f} 
% % \right)
% % \right|
% % \\
% % &\leq
% % \E \left(
% %   \left|
% %   \prod_{ef \in \gamma} \bar{B}_{e f} 
% %   \right|
% % \right)  \quad \mbox{Jensen's inequality}
% % \\
% % &=
% % \E \left(
% %     | X_\gamma |
% % \right)
% % \end{align*}

% Let $E(\gamma) = \{ e_1, e_2, \ldots, e_\nedge \}$ be the set of 
% undirected edges visited by $\gamma$,
% and let $H_t$ be the graph induced by the set of undirected edges 
% $\{e_i\}_{i=1}^t$.
% To get an adequate upper bound,
% we order the edges so that we have a maximal number of good edges
% with two-path multiplicity one.
% Let $F$ be the set of undirected edges whose oriented counterparts,
% in either direction, 
% participate in a two-path which occurs only once in $\gamma$.
% Then,
% \begin{align*}
% F 
% % &= 
% % \bigcup_{t=1}^\nedge
% % \{ e_t : 
% % 	  m_{\vec{e}_t f} + m_{f^{-1} (\vec{e}_t)^{-1}} =1 
% %       \mbox{ or } 
% % 	  m_{g \vec{e}_t} + m_{ (\vec{e}_t)^{-1} g^{-1}} = 1
% %       \mbox{ for any $f,g \in \vec{E}(\gamma)$ }
% % \} 
% % \\
% &= 
% \{
% e \in E(\gamma):
% m_{e f} + m_{f^{-1} e^{-1}} = 1 
% \mbox{ for some two-path $e f \in \gamma$ }
% \} .
% \end{align*}
% We assume the edges have already been ordered
% and will revisit how to pick the ordering $\pi$ at the end of the proof.

% Now denote $\vec{A}_t \subset \vec{E}(\gamma \cup \gamma^{-1})$ 
% to be the set of {\it oriented} edges
% \[
% \vec{A}_t=\{\vec{e}_1,(\vec{e}_1)^{-1}, \dots, 
% \vec{e}_t, (\vec{e}_t)^{-1}\}
% \]
% containing the first $t$ edges with both possible orientations,
% irregardless of whether that orientation occurs in $\gamma$.
% For simplicity and without any loss of generality,
% we make the convention that $\vec{e}_t$ 
% goes from set $V_1$ to $V_2$.
% Define $\mathcal{F}_t$ as the sigma algebra generated by $\vec{A}_t$
% for $t = 1, 2, \ldots, \nedge$.
% Then $\mathcal{F}_1\subset \mathcal{F}_2\subset \cdots \subset \mathcal{F}_{\nedge}$,
% so $\{ \mathcal{F}_t \}$ is a filtration. 


% Recall the definition of $X_{\gamma}$ from Eqn.~\eqref{def:X_gamma}. 
% We define the random variables $\{X_t\}_{2\leq t\leq \nedge}$ as
% \[
% X_t = 
% \prod_{ef}
% % _{ 
% % \scriptsize
% % \begin{array}{c}
% % ef \in \gamma \\
% % e, f \in \vec{A}_t
% % \end{array}
% % } 
% \left| \bar{B}_{ef} \right|^{m_{ef}} ,
% \]
% where the product is over all 
% $e,f \in \vec{A}_t$, where $ef$ is a subpath of $\gamma$, 
% and 
% exactly one of $\{e, f\}$
% is in the set 
% $\vec{A}_t \setminus \vec{A}_{t-1} = \{\vec{e}_t,(\vec{e}_t)^{-1}\}$. 
% Thus, $X_t$ is the product of all factors in $X_{\gamma}$ 
% which involve $e_t$ at its first appearance under the ordering. 
% Note that the following definition is equivalent:
% \begin{equation}
% \label{eq:Xt_mult}
% X_t = 
% \prod_{e, f_0, f_1, f_2, f_3}
% \left|\bar{B}_{e f_0} \right|^{m_{e f_0}}
% \left|\bar{B}_{e^{-1} f_1} \right|^{m_{e^{-1} f_1}}
% \left|\bar{B}_{f_2 e} \right|^{m_{f_2 e}}
% \left|\bar{B}_{f_3 e^{-1}} \right|^{m_{f_3 e^{-1}}},
% \end{equation}
% where $e, e^{-1} \in \vec{A}_t \setminus \vec{A}_{t-1}$, 
% and $f_0, f_1, f_2, f_3 \in \vec{A}_{t-1}$.
% It is not hard to see that \textcolor{green}{no need for absolute value since the $X_t$ are now positive. GB}
% $
% |X_{\gamma}| = \prod_{t=2}^\nedge X_t .
% $
% Defining
% $
% Y_t=\prod_{j=2}^t X_j
% $ 
% for $2 \leq t \leq \nedge$,
% we see that 
% $|X_\gamma| = Y_\nedge$
% and
% $Y_t = X_t Y_{t-1}$ for $2 < t \leq \nedge$.
% Also, every $Y_t$
% is $\mathcal{F}_t-$measurable.
% With these ingredients, 
% by the law of total probability we have
% the following tower of expectation:
% \begin{align}
% \label{eq:tower_expectation}
% \E(|X_{\gamma}|) =
% \E(\E( Y_\nedge | \mathcal{F}_{\nedge-1} )) =
% \E(\E( Y_{\nedge-1} X_\nedge | \mathcal{F}_{\nedge-1} )) =
% \E(Y_{\nedge-1} \E( X_{\nedge} | \mathcal{F}_{\nedge-1} )).
% \end{align}

% We focus on the term $\E(X_{\nedge}|\mathcal{F}_{\nedge-1}))$
% in Eqn.~\ref{eq:tower_expectation}. 
% For oriented edges $e, f_0, f_1, f_2, f_3$ as above,
% we have that the conditional random variables
% \[
% (B_{e f_0} | f_0) 
% = (B_{e^{-1} f_1} | f_1)
% = (B_{f_2 e} | f_2) 
% = (B_{f_3 e^{-1}} | f_3) ,
% \]
% where we assume that the orientation of the edges
% is such that these entries are not zero. 
% These equalities say that, 
% under the event that $f_0, f_1, f_2, f_3$ are edges of the graph, 
% the presence of edge $e$ or $e^{-1}$ attached to any of them
% is equivalent, 
% making these random variables identical.
% Note that this involves the non-backtracking matrix $B$, 
% not its ``centered'' version $\bar{B}$.
% These equivalences allow us to combine the multiplicities in Eqn.~\ref{eq:Xt_mult}.

% Let $\{Z_t\}_{2\leq t\leq \nedge}$ be independent random variables with distribution
% \[
% Z_t 
% \eqd 
% ( B_{\vec{e}_t f} \, |f ) ,
% \]   
% where $f\in \vec{A}_{t-1}$, 
% $\vec{e}_t \in \vec{A}_t \setminus \vec{A}_{t-1}$,
% and $\vec{e}_t f$ is a subpath of $\gamma$.
% Since our convention is that $\vec{e}_\nedge$ terminates in set $V_2$, 
% we have
% \begin{align}
% \label{exp of X_eta}
% \E(X_{\nedge}|\mathcal{F}_{\nedge-1})
% =
% \E 
% \left( 
% \left| Z_\nedge - \frac{d_1-1}{m} \right|^a
% \left| Z_\nedge-\frac{d_2-1}{n} \right|^b
% \right),
% \end{align}
% where $a$ is the number of times the new edge $\vec{e}_\nedge$
% appears as $B_{f_2 \vec{e}_\nedge}$ and $B_{(\vec{e}_\nedge)^{-1} f_1}$ in $X_t$,
% and $b$ is the number of times it appears as 
% $B_{\vec{e}_\nedge f_0}$ and $B_{f_3 (\vec{e}_\nedge)^{-1}}$.
% That is,
% \[
% a = \sum_{f_1,f_2 \in \vec{A}_{\nedge-1}} m_{(\vec{e}_\nedge)^{-1} f_1} + m_{f_2 \vec{e}_\nedge}
% \]
% and
% \[
% b = \sum_{f_0,f_3 \in \vec{A}_{\nedge-1}} m_{\vec{e}_\nedge f_0} + m_{f_3 (\vec{e}_\nedge)^{-1}} .
% \]

% \todo{I feel like we should make explicit what $r$ is when we apply 
% Lemma 11 here}\textcolor{green}{Agree!}

% To evaluate Eqn.~\eqref{exp of X_eta}, 
% we will use the fact that $Z_\nedge$ is a Bernoulli random variable
% and $0 < (d_2-1)/n \leq (d_1-1)/m < 1$.
% Therefore, the following hold deterministically:
% \begin{equation*}
%   \left| Z_\nedge - \frac{d_1-1}{m} \right| \leq 1 
%   \quad \mbox{and} \quad
%   \left| Z_\nedge - \frac{d_2-1}{n} \right| \leq 1. 
% \end{equation*}
% We have several cases depending on the values of $a$ and $b$.
% \begin{enumerate}[{Case} 1:]
% \item If $a\geq 2$, regardless of $b$, then   
%   \[
%   \left| Z_\nedge - \frac{d_1-1}{m} \right|^a
%   \left| Z_\nedge-\frac{d_2-1}{n} \right|^b 
%   \leq 
%   \left| Z_\nedge - \frac{d_1-1}{m} \right|^2.
%   \]
% By Lemmas \ref{bernoulli_bound} and \ref{one_edge}, 
%   \[
%   \E \left| Z_\nedge - \frac{d_1-1}{m} \right|^2
%   \leq 
%   \frac{d_1-1}{m} + O\left(\frac{\nedge}{n^2}\right). 
%   \] 
%   When $b \geq 2$, regardless of $a$, the bound is analogous with 
%   $(d_2 - 1)/n$ instead.
% %\item $s\geq 2$. Is analogous.
% \item If $a=b=1$, recall that $Z_\nedge$ is a Bernoulli
%     random variable with some success probability 
%     $q$.
%     Letting 
%     $p_1 = \frac{d_1-1}{m}$
%     and 
%     $p_2 = \frac{d_2-1}{n}$,
%     we have that
%     $q = p_1 + O(\frac{\nedge}{n^2})$
%     if the new edge attaches in set $V_1$,
%     and $q = p_2 + O(\frac{\nedge}{n^2})$
%     if the new edge attaches in set $V_2$.
%   \begin{align*}
%   \E \left( 
%   	\left| Z_\nedge - \frac{d_1-1}{m} \right| 
%   	\left| Z_\nedge - \frac{d_2-1}{n} \right|
%     \right)
%     &= 
%     (1-p_1)(1-p_2) q + p_1 p_2 (1-q) \\
%     &= q + O \left( \frac{\nedge}{n^2} \right) \\
%     &=
%     \left\{
%     \begin{array}{cc}	
%         \frac{d_1-1}{m} + O\left(\frac{\nedge}{n^2}\right), & 
%         	\mbox{$e_\nedge$ attaches in $V_1$} \\
%      	\frac{d_2-1}{n} + O\left(\frac{\nedge}{n^2}\right), & 
%         	\mbox{$e_\nedge$ attaches in $V_2$} 
%      	\end{array}
%      \right. .
%   \end{align*}
% \item If $a=0$ and $b=1$,
% 	by \todo{new Lemma}, 
%     for Eqn.~\eqref{exp of X_eta} we then have
% 	\[
% 	\E \left| Z_\nedge - \frac{d_2-1}{n} \right|
%      \leq
%      2\left( \frac{d_2-1}{n} \right) + O\left(\frac{\nedge}{n^2}\right).
% %      \left\{
% % 		\begin{array}{cc}	
% %         \frac{d_2-1}{n} + O\left(\frac{\nedge}{n^2}\right), & 
% %         	\mbox{$e_\nedge$ is not a good edge} \\
% %      	O\left(\frac{\nedge}{n^2}\right), &
% %         	\mbox{$e_\nedge$ is a good edge}
% %         \end{array}
% %      \right.
% 	\]
% \item If $b=0$ and $a=1$,
% 	by \todo{new Lemma}, 
%     for Eqn.~\eqref{exp of X_eta} we then have
%     \[
% 	\E \left| Z_\nedge - \frac{d_1-1}{m} \right|
%      \leq
%      2 \left( \frac{d_1-1}{m} \right) + O\left(\frac{\nedge}{n^2}\right) .
% %      \left\{
% % 		\begin{array}{cc}	
% %         \frac{d_1-1}{m} + O\left(\frac{\nedge}{n^2}\right), & 
% %         	\mbox{$e_\nedge$ is not a good edge} \\
% %      	O\left(\frac{\nedge}{n^2}\right), &
% %         	\mbox{$e_\nedge$ is a good edge}
% %         \end{array}
% %      \right.
% 	\]
% \end{enumerate} 

% \todo{With the updates to above cases 1--4, would need to rewrite the rest here}

% We summarize these cases, noting that
% adding edge $e_\nedge$ to $H_{\nedge-1}$
% results in a different upper bound for 
% $\E(X_\nedge | \mathcal{F}_{\nedge-1})$
% depending on: 
% (1) the multiplicity of the two-paths in which it appears,
% (2) whether it connects to a vertex in $V_1$ or $V_2$
% when it attaches to $H_{\nedge-1}$, and
% (3) whether it is a good edge.
% Note that cases 1--3 are ambiguous about what bipartite
% set $e_\nedge$ connects to in the graph $H_{\nedge-1}$.
% If both endpoints of $e_\nedge$ connect to edges in $H_{\nedge-1}$, 
% then it is not a good edge,
% and Lemma~\ref{one_edge} part (iii) says either degree bound is applicable.

% Using Eqn.~\ref{exp of X_eta} and 
% summarizing all of the cases enumerated above,
% Eqn.~\eqref{eq:tower_expectation} becomes
% \begin{align}
% \E(Y_{\nedge}) &= 
% \E(Y_{\nedge-1} \E(X_\nedge|\mathcal{F}_{\nedge-1})) \nonumber \\
% &\leq
% \E( Y_{\nedge-1} ) \, \E(X_\nedge|\mathcal{F}_{\nedge-1}))
% \nonumber \\
% &\leq 
% \E(Y_{\nedge -1 } )
% \left\{
% \begin{array}{cc}
% \frac{ d_1 - 1 }{m} + O \left( \frac{\nedge}{n^2}
% \right), 
% & \mbox{$e_\nedge$ connects to $V_1$ or both and $a+b>1$} \\
% \frac{ d_2 - 1 }{n} + O \left( \frac{\nedge}{n^2}
% \right), 
% & \mbox{$e_\nedge$ connects to $V_2$ or both and $a+b>1$} \\
% O \left( \frac{\nedge}{n^2} \right), 
% & \mbox{$e_\nedge$ is good and $a+b=1$ } \\
% \frac{ d_2 }{n} + O \left( \frac{\nedge}{n^2}
% \right), 
% & \mbox{$e_\nedge$ is isolated from $H_{\nedge-1}$ }.
% \end{array}
% \right.
% \label{eq:case_cond_expectation}
% \end{align}
% The first inequality follows from the fact that
% $X_\nedge \geq 0$ and $Y_{\nedge -1 } \geq 0$.
% In Eqn.~\ref{eq:case_cond_expectation}, we again see that 
% if $e_\nedge$ connects two vertices already in $H_{\nedge-1}$
% we are free to choose which bound.

% We apply the same argument to 
% $\E(Y_{\nedge-1}) = \E( Y_{\nedge-2} \E(X_{\nedge-1} | \mathcal{F}_{\nedge-2}) )$, 
% and continue this process.
% After a total $\nedge-2$ times, we get to 
% $
% \E (Y_2) = \E (X_2) = \E ( \E( X_2 | \mathcal{F}_1 ) )
% $.
% The conditional expectation 
% $ \E( X_2 | \mathcal{F}_1 )$
% is just another term
% like in Eqn.~\eqref{eq:case_cond_expectation},
% and then there is a final term that comes from the
% first, isolated edge.

% After all of these iterations, 
% we end up with a bound of the form:
% \begin{align}
% \E(Y_\nedge)
% & \leq
% \left( \frac{d_1-1}{m} \right)^r
% \left( \frac{d_2-1}{n} \right)^{\nedge-r-s}
% \left( \frac{d_2}{n} \right)^{s}
% \left( \frac{\nedge}{n} \right)^\omega
% \left(1 + o(1) \right) 
% %\label{eqn:bound_strong}
% \nonumber 
% \\
% & = 
% \left( \frac{d_1-1}{m} \right)^r
% \left( \frac{d_2-1}{n} \right)^{\nedge-r}
% \left( \frac{d_2}{d_2 - 1} \right)^{s}
% \left( \frac{\nedge}{n} \right)^\omega
% \left(1 + o(1) \right) .
% \label{eqn:bound_strong}
% % \\
% % & \leq 
% % \frac{d^\nedge}{m^r n^{\nedge-r}} 
% % \left( \frac{\nedge}{n} \right)^\omega
% % \left(1 + o(1) \right) .
% % \label{eqn:bound_weak}
% % \\
% % & = 
% % \frac{D_\gamma}{m^Q n^{\nedge-Q}}
% % \left( \frac{\nedge}{n} \right)^\omega
% % \left(1 + o(1) \right) .
% \end{align}
% Here, $r$ is the number of times 
% the added edge attaches only in set $V_1$
% (or in both, 
% in which case we choose the bound which balances
% $r$ and $\nvright$),
% whereas $s$ is the number of isolated edges which are added.
% % The powers of $(d_1 - 1)$ and $(d_2 - 1)$ are counted in 
% % $D_\gamma = (d_1 -1 )^Q (d_2 - 1)^{\nedge-Q}$.
% The exponent $\omega$ counts the number of times we add 
% a good edge with $a+b = 1$,
% which gives an extra factor of $\nedge/n$ in the expectation.
% For many good edges guaranteed by our ordering, 
% we add an isolated edge first,
% but at most $6d^2 \omega$ isolated edges are added in this way.
% \todo{This isn't obvious to me... we should justify}

% Now we consider two cases, depending on whether $F$ is empty or not.
% If all edges appear at least twice, then $F$ is empty
% and $\omega = 0$.
% Since the ordering $\pi$ is arbitrary, 
% consider the ordering induced by the circuit $\gamma$, 
% where we add the edges which create cycles last.
% Thus, only the first edge is isolated and $s = 1$.
% Furthermore, each additional edge contributes factors of
% $(d_1-1)/m$ and $(d_2-1)/n$ in alternation,
% depending on the set of the first vertex.
% We add the non-tree edges last, and since these connect to both bipartite sets, 
% we can choose either the $(d_1-1)/m$ or $(d_2-1)/n$ bound.
% We can {\em always} pick these bounds in such a way that $r=\nvright$,
% because $\gamma$ is a circuit on a bipartite graph. 
% Thus, in this case, we employ Eqn.~\ref{eqn:bound_strong} to find that
% \begin{align*}
% \E(X_{\gamma})
% & \leq
% \left( \frac{d_1-1}{m} \right)^\nvright
% \left( \frac{d_2-1}{n} \right)^{\nedge-\nvright}
% \left( \frac{d_2}{d_2-1} \right)
% \left(1 + o(1) \right) .
% \end{align*}
% The factor of $\frac{d_2}{d_2-1}$ comes from the
% single $s=1$ isolated vertex.
% Assuming $d_2 \leq d_1$, 
% this holds regardless of whether we use $d_1/m$ or $d_2/n$ for the
% isolated, starting edge.

% Otherwise, $F$ is nonempty.
% First, we order the edges in $F$ following the algorithm
% in Lemma~\ref{lem:good_ordering}, producing an ordering $\tilde{\pi}$.
% The actual ordering $\pi$ is set by using
% $\pi(i) = \tilde{\pi}(i)$ where it is defined
% and arbitrarily extending $\pi$ to the rest of the edges.

% The order that we add the edges allows us to control $r$.
% First of all, consider the edges in $F$ that are ordered 
% according to the algorithm in Lemma~\ref{lem:good_ordering}.
% When these are added, we can choose which bound to employ so that
% all $\nvright$ vertices in $V_2$ contribute to $r$:
% Isolated edges or edges which connect to both bipartite sets allow us freedom 
% to choose $(d_2-1)/n$ or $(d_1-1)/m$.
% The ordering process in Lemma~\ref{lem:good_ordering} 
% stops at some time $t = T$.
% For $t > T$, pick some $u \in V_1 \setminus V(H_{t-1})$ 
% and an edge $e$ incident to $u$.
% If this edge $e = (u,v)$ connects to a vertex $v$ already in the current $V(H_{t-1})$,
% then $v$ already has a factor in $r$.
% If $v \not \in V(H_{t-1})$, it is a currently isolated edge, 
% and we take the bound that allows us to increase $r$.
% Once every $u \in V_1$ is contained in $H_{t-1}$, 
% any new edge will connect to just $V_1$ or both, and we can add one to
% $r$ for each $v \in V_2$ that needs it.
% Once $r=\nvright$, every vertex has been discovered, and thus any edges left over
% will connect to both sets $V_1$ and $V_2$.
% In this case, we use the bound that does not increase $r$.
% Thus, we can insure that $r=\nvright$.

% Let $p_2(F)$ denote the number of two-paths in $F$,
% then 
% $p_2(F) \geq 
% \sum_{ef \in \gamma} 1 \{m_{ef} + m_{f^{-1} e^{-1}} = 1 \}$.  
% By Lemma~\ref{lem:good_ordering}, 
% we construct an ordering $\pi$
% which guarantees that 
% $\left\lceil p_2(F) / (6d^2) \right\rceil$
% edges are good,
% and these edges must also have $a+b = 1$.
% % For a guaranteed good edge $e_{2t} \in F$,
% % all of the edges in
% % $N(\{ e_{2t} \} ) \setminus E(H_{2t-1})$ appear later.
% %Since $e_{2t}$ occurs in a two-path of multiplicity one
% %in $H_{2t}$,
% %by construction, then $a+b=1$.
% So this gives a power of at least
% \[
%  \omega =
%  \left\lceil 
%  \frac{\sum_{ef \in \gamma} 
%  1 \{m_{ef} + m_{f^{-1} e^{-1}} = 1 \}}{6d^2} 
%  \right\rceil .
% \]

% This shows that
% \begin{align*}
% \mathbb{E} \left( |X_{\gamma}| \right)
% &\leq 
% \frac{ (d_1-1)^{\nvright}(d_2-1)^{\nedge - \nvright} }{ m^\nvright n^{\nedge - \nvright} }  
% \left( \frac{d_2}{d_2 - 1} \right)^s
% \left( \frac{\nedge}{n} \right)^{ \omega }
% (1 + o(1)) & & \\
% &=
% \frac{ ((d_1 - 1) (d_2 - 1))^{\nedge/2} }{ m^\nvright n^{\nedge - \nvright} }  
% \left( \frac{d_1 - 1}{d_2 - 1} \right)^{\nvright - \nedge/2}
% %\left( d - 1 \right)^{\nedge - 2 k \ell}
% \left( \frac{d_2}{d_2 - 1} \right)^s
% \left( \frac{\nedge}{n} \right)^{ \omega }
% (1 + o(1)) & & \\
% &\leq
% \frac{ ((d_1 - 1) (d_2 - 1))^{\nedge /2} }{ m^\nvright n^{\nedge - \nvright} }  
% \left( \frac{d_1 - 1}{d_2 - 1} \right)^{\nvright - \nedge/2}
% %\left( d - 1 \right)^{\nedge - 2 k \ell}
% \left( \frac{c \nedge}{n} \right)^{ \omega }
% (1 + o(1)) , & & 
% (s \leq 6d^2 \omega \mbox{, Ioana's bound})
% \end{align*}
% as claimed.
% \end{proof}

